\documentclass[a4paper]{article}
\usepackage{ctex}
\ctexset{proofname=\heiti{证明}}
\usepackage{amsmath,amssymb,amsthm}
\usepackage{mathtools}
\usepackage{bm}
\usepackage{enumitem}
\usepackage{booktabs}
\usepackage{array}
\usepackage{graphicx}

% 若本地有课程模板 iidef.sty，则优先使用；
% 若没有，则使用备用定义，保证文件尽量可以直接编译。
\IfFileExists{iidef.sty}{
    \usepackage[thehwcnt = 6]{iidef}
    \thecourseinstitute{清华大学}
    \thecoursename{人工智能原理}
    \theterm{2026年春季学期}
    \hwname{作业}
    \slname{\heiti{解}}
}{
    \makeatletter
    \newcommand{\thecourseinstitute}[1]{\def\@courseinstitute{#1}}
    \newcommand{\thecoursename}[1]{\def\@coursename{#1}}
    \newcommand{\theterm}[1]{\def\@term{#1}}
    \newcommand{\hwname}[1]{\def\@hwname{#1}}
    \newcommand{\slname}[1]{\def\@slname{#1}}
    \thecourseinstitute{清华大学}
    \thecoursename{人工智能原理}
    \theterm{2026年春季学期}
    \hwname{作业}
    \slname{\heiti{解}}
    \newcommand{\courseheader}{
        \begin{center}
            {\Large \heiti \@courseinstitute \quad \@coursename}\\[0.5em]
            {\large \@term \quad \@hwname 6 \quad \@slname}
        \end{center}
        \vspace{1em}
    }
    \newcommand{\name}[1]{\noindent #1 \par\vspace{1em}}
    \newenvironment{solution}{\par\noindent\textbf{解：}\par}{\par}
    \makeatother
}

\begin{document}
\courseheader
\name{姓名：杜一郎\qquad 学号：2023011618}

\begin{enumerate}
\setlength{\itemsep}{3\parskip}

\item CNN 输出尺寸与参数量
\begin{solution}
\textbf{(1) 计算卷积层 1 和卷积层 2 输出特征图的尺寸}

对卷积层或池化层，若输入空间尺寸为 $H_{\mathrm{in}}\times W_{\mathrm{in}}$，卷积核或窗口大小为 $K$，填充为 $P$，步长为 $S$，则输出空间尺寸为
\[
H_{\mathrm{out}}
=
\left\lfloor
\frac{H_{\mathrm{in}}+2P-K}{S}
\right\rfloor+1,
\qquad
W_{\mathrm{out}}
=
\left\lfloor
\frac{W_{\mathrm{in}}+2P-K}{S}
\right\rfloor+1.
\]

卷积层 1 的输入为 $32\times 32\times 3$，卷积核空间大小为 $5\times 5$，步长 $S=1$，填充 $P=2$，因此
\[
H_{\mathrm{out}}
=
\left\lfloor
\frac{32+2\times 2-5}{1}
\right\rfloor+1
=
32,
\]
\[
W_{\mathrm{out}}
=
\left\lfloor
\frac{32+2\times 2-5}{1}
\right\rfloor+1
=
32.
\]
卷积核数量为 $16$，所以卷积层 1 的输出特征图尺寸为
\[
32\times 32\times 16.
\]

接着经过最大池化层。池化层输入为 $32\times 32\times 16$，窗口大小 $2\times 2$，步长 $2$，不额外填充，因此
\[
H_{\mathrm{out}}
=
\left\lfloor
\frac{32-2}{2}
\right\rfloor+1
=
16,
\]
\[
W_{\mathrm{out}}
=
\left\lfloor
\frac{32-2}{2}
\right\rfloor+1
=
16.
\]
池化层不改变通道数，因此池化层输出为
\[
16\times 16\times 16=4096.
\]

卷积层 2 的输入为 $16\times 16\times 16$，卷积核空间大小为 $3\times 3$，步长 $S=1$，填充 $P=1$，所以
\[
H_{\mathrm{out}}
=
\left\lfloor
\frac{16+2\times 1-3}{1}
\right\rfloor+1
=
16,
\]
\[
W_{\mathrm{out}}
=
\left\lfloor
\frac{16+2\times 1-3}{1}
\right\rfloor+1
=
16.
\]
卷积层 2 的卷积核数量为 $32$，因此卷积层 2 的输出特征图尺寸为
\[
16\times 16\times 32=8192.
\]

\textbf{(2) 计算卷积层 1 的可训练参数数量，并与直接展平输入连接到全连接层的参数数量进行比较}

卷积层 1 中每个卷积核大小为 $5\times 5\times 3$，因此每个卷积核的权重参数数量为
\[
5\times 5\times 3=75.
\]
卷积层 1 共有 $16$ 个卷积核，所以权重参数总数为
\[
75\times 16=1200.
\]
若每个输出通道对应一个偏置参数，则偏置参数数量为 16.

因此卷积层 1 的可训练参数总数为
\[
1200+16=1216.
\]

\textbf{(3) 直接展平输入并连接到 128 个神经元的全连接层}

原始输入图像大小为 $32\times 32\times 3$，展平后的一维向量维度为
\[
32\times 32\times 3=3072.
\]
如果将该向量连接到一个具有 $128$ 个神经元的全连接层，则权重参数数量为
\[
3072\times 128=393216.
\]
每个神经元有一个偏置参数，因此偏置参数数量为128.

于是该全连接层的可训练参数总数为
\[
393216+128=393344.
\]

与卷积层 1 的参数数量 $1216$ 相比，全连接层的参数数量明显更多。二者的比值为
\[
\frac{393344}{1216}\approx 323.47.
\]
可见卷积层通过局部连接和权重共享，大幅减少了模型参数量。
\end{solution}

\item 感受野与空洞卷积

\begin{solution}
考虑计算下面两个量：
\begin{itemize}
    \item $r_l$：第 $l$ 层一个神经元在原始输入图像上的感受野大小；
    \item $j_l$：第 $l$ 层相邻两个神经元在原始输入图像上对应位置的间隔。
\end{itemize}

初始化输入层为
\[
r_0=1,\qquad j_0=1.
\]
对于普通卷积或池化层，若卷积核或窗口大小为 $k_l$，步长为 $s_l$，则
\[
r_l=r_{l-1}+(k_l-1)j_{l-1},
\]
\[
j_l=j_{l-1}s_l.
\]
填充会影响特征图的边界位置和尺寸，但不改变单个神经元理论感受野的大小。

\textbf{(1) 计算 C2 输出特征图上一个神经元在原始输入图像上的感受野大小}

首先经过卷积层 C1。C1 的卷积核大小为 $5\times 5$，步长为 $1$，所以
\[
r_{\mathrm{C1}}
=
r_0+(5-1)j_0
=
1+4\times 1
=
5,
\]
\[
j_{\mathrm{C1}}
=
j_0\times 1
=
1.
\]
因此，C1 输出层一个神经元在原始图像上的感受野为 $5\times 5$。

接着经过池化层 P1。P1 的窗口大小为 $2\times 2$，步长为 $2$，所以
\[
r_{\mathrm{P1}}
=
r_{\mathrm{C1}}+(2-1)j_{\mathrm{C1}}
=
5+1\times 1
=
6,
\]
\[
j_{\mathrm{P1}}
=
j_{\mathrm{C1}}\times 2
=
2.
\]

最后经过卷积层 C2。C2 的卷积核大小为 $3\times 3$，步长为 $1$，所以
\[
r_{\mathrm{C2}}
=
r_{\mathrm{P1}}+(3-1)j_{\mathrm{P1}}
=
6+2\times 2
=
10,
\]
\[
j_{\mathrm{C2}}
=
j_{\mathrm{P1}}\times 1
=
2.
\]
因此，C2 输出特征图上的一个神经元在原始输入图像上的感受野大小为
\[
10\times 10=100.
\]

\textbf{(2) 若去掉池化层 P1，C2 输出神经元的感受野大小}

如果去掉池化层 P1，则网络结构变为 C1 直接连接 C2.

C1 部分不变：
\[
r_{\mathrm{C1}}=5,\qquad j_{\mathrm{C1}}=1.
\]
然后直接经过 C2。C2 的卷积核大小为 $3\times 3$，步长为 $1$，所以
\[
r_{\mathrm{C2}}
=
r_{\mathrm{C1}}+(3-1)j_{\mathrm{C1}}
=
5+2\times 1
=
7,
\]
\[
j_{\mathrm{C2}}
=
j_{\mathrm{C1}}\times 1
=
1.
\]
因此去掉池化层后，C2 输出神经元的感受野大小为
\[
7\times 7=49.
\]

可以看到，有池化层时感受野为 $10\times 10$，没有池化层时为 $7\times 7$。池化层虽然本身没有可训练参数，但它有窗口大小和步长。窗口大小会直接扩大当前层的感受野，步长大于 $1$ 则会增大后续层相邻神经元在原始图像中的间隔，使后续卷积层每增加一层时感受野增长得更快。因此，池化层通常会加速感受野的扩大。

\textbf{(3) 将 C2 改为空洞率为 2 的空洞卷积时的感受野大小}

空洞卷积的有效卷积核大小为
\[
k_{\mathrm{eff}}
=
k+(k-1)(d-1),
\]
其中 $k$ 是原始卷积核大小，$d$ 是空洞率。

本题中 C2 的卷积核大小仍为 $3\times 3$，空洞率为 $d=2$，因此一维方向上的有效卷积核大小为
\[
k_{\mathrm{eff}}
=
3+(3-1)(2-1)
=
5.
\]
也就是说，一个 $3\times 3$、空洞率为 $2$ 的空洞卷积在空间覆盖范围上等价于一个 $5\times 5$ 的采样范围，但实际参数量仍然是 $3\times 3$ 个卷积核位置对应的参数。

保持原结构，即仍然有 C1 和 P1。前两层的感受野计算结果仍为
\[
r_{\mathrm{P1}}=6,\qquad j_{\mathrm{P1}}=2.
\]
现在 C2 使用有效卷积核大小 $k_{\mathrm{eff}}=5$，步长仍为 $1$，所以
\[
r_{\mathrm{C2}}
=
r_{\mathrm{P1}}+(k_{\mathrm{eff}}-1)j_{\mathrm{P1}}
=
6+(5-1)\times 2
=
14.
\]
因此此时 C2 输出神经元在原始输入图像上的感受野大小为
\[
14\times 14=196.
\]

相比普通 $3\times 3$ 卷积得到的 $10\times 10$ 感受野，空洞率为 $2$ 的空洞卷积可以在不增加卷积核参数数量的情况下，把该层输出神经元的感受野扩大到 $14\times 14$。
\end{solution}

\end{enumerate}

\end{document}
